\documentclass[11pt]{article}
\usepackage[margin=1in]{geometry}
\usepackage{amsmath,amssymb}
\usepackage{graphicx}
\usepackage{booktabs}
\usepackage{hyperref}
\usepackage{microtype}

\title{\textbf{Paper 42: Three Anomalies Closed ---
Spontaneous Copy-Forward Maintenance,\\
the C$_\mathrm{half}$ Surge Mechanism,
and Checkpoint Phase Aliasing}}
\author{VCSM Theoretical Series}
\date{2026-03-10}

\begin{document}
\maketitle

\begin{abstract}
Papers~38--41 established the core single-region VCML theory but left three
quantitative loose ends.
Anomaly~1: in the V89 location-crystal condition, some seeds showed
zone-structure ratio $> 1$ (structure growing without new wave input).
Anomaly~2: in Paper~40, the half-width boundary shift ($\Delta=5$) caused
sg4n to reach $0.80$ at $T=4000$, far exceeding the C$_\mathrm{ref}$ peak
of $0.38$.
Anomaly~3: in Paper~41, sg4n$_B > $ sg4n$_A$ at $T_\mathrm{switch}=100$
steps but sg4n$_A > $ sg4n$_B$ at $T_\mathrm{switch}=50$ steps --- a
reversal contradicting the first-mover advantage.
We resolve all three.
Anomaly~1 (\textit{Exp~A -- wave ablation}): when waves stop at the
commitment peak ($T=2000$), structure decays with ratio $0.664$, \emph{above}
the field-decay-only prediction of $0.549$.
Spontaneous collapses (from ageing and instability noise) continue to trigger
copy-forward births even without active wave input, partially offsetting
$\mathrm{FIELD\_DECAY}$ and explaining why crystals with rare turnover can
maintain or even grow sg4n.
Anomaly~2 (\textit{Exp~B -- surge anatomy}):
the C$_\mathrm{half}$ surge persists to $T=6000$, where sg4n$=0.21$
versus C$_\mathrm{ref}=0.07$ (3.1$\times$).
The mechanism is the same as Paper~41's copy-forward refreshing:
zone-boundary mismatches after the shift drive sustained structured collapses,
continuously replenishing fieldM.
Anomaly~3 (\textit{Exp~C -- phase aliasing}):
the 300-step checkpoint interval is an even multiple of $T_\mathrm{switch}=50$
($300/50=6$), causing all 20 CPS measurements to fall in task~B epochs.
For $T_\mathrm{switch}=100$ ($300/100=3$, odd), the last three checkpoints
fall 2:1 in task~B, inflating sg4n$_B$.
Epoch-end measurements (unbiased) show sg4n$_A \geq $ sg4n$_B$ for both
switching periods, restoring the expected first-mover ordering.
Single-region VCML theory is closed with no remaining anomalies.
\end{abstract}

\section{Introduction}

Papers~38--41 completed the characterisation of single-region VCML dynamics:
metastability half-life $\approx 800$ steps (Paper~38); permutation-transparent
continual learning (Paper~40); copy-forward refreshing as an anti-decay
mechanism (Paper~41).
Three quantitative results resisted clean interpretation:
\begin{enumerate}
  \item \textbf{V89 location-crystal ratio $>1$:}
        In the C$_\mathrm{loc}$ crystal condition (cells constrained to fixed
        positions, births at the same location), some seeds showed
        sg4$_\mathrm{final}$/sg4$_\mathrm{initial} > 1$.
        Since $\mathrm{FIELD\_DECAY}=0.9997$ decays fieldM every step,
        growing sg4n without fresh wave input requires an alternative
        replenishment mechanism.
  \item \textbf{C$_\mathrm{half}$ late surge (Paper~40):}
        After a half-width zone shift ($\Delta=5$), sg4n reached $0.80$ at
        $T=4000$, more than twice the C$_\mathrm{ref}$ peak ($0.38$).
        The Paper~41 copy-forward mechanism was proposed as the cause, but
        was not directly measured.
  \item \textbf{$T_\mathrm{switch}=50$ vs.\ 100 reversal (Paper~41):}
        At $T_\mathrm{switch}=50$, sg4n$_A > $ sg4n$_B$.
        At $T_\mathrm{switch}=100$, sg4n$_B > $ sg4n$_A$.
        This ordering flip contradicts the first-mover prediction and any
        monotone relationship between switching period and task advantage.
\end{enumerate}

We address each anomaly with a targeted experiment, then unify the findings.

\section{Methods}

Standard VCML parameters: $H=40$, $\mathrm{HALF}=40$, $N_\mathrm{zones}=4$,
$z_w=10$, $\mathrm{WR}=4.8$, $\mathrm{FA}=0.16$, 5 seeds per condition
(Exp~A and B), 10 seeds (Exp~C).

\paragraph{Exp A: Wave ablation (Anomaly~1).}
Three conditions, $T=4000$, checkpoints every 200 steps:
\begin{itemize}
  \item \textbf{C$_\mathrm{ref}$}: standard waves throughout ($\mathrm{WR}=4.8$).
  \item \textbf{C$_\mathrm{wavestop}$}: normal waves until $T_\mathrm{stop}=2000$
        (the commitment peak), then $\mathrm{WR}=0$ and wave queue cleared.
  \item \textbf{C$_\mathrm{nowaves}$}: $\mathrm{WR}=0$ from $t=0$ (never any waves).
\end{itemize}
Metric: sg4n and cumulative collapse count at each checkpoint.

\paragraph{Exp B: C$_\mathrm{half}$ surge anatomy (Anomaly~2).}
Two conditions, $T=6000$, checkpoints every 300 steps:
\begin{itemize}
  \item \textbf{C$_\mathrm{ref}$}: $\Delta=0$ throughout.
  \item \textbf{C$_\mathrm{half}$}: $\Delta=0$ until $T_\mathrm{shift}=1500$,
        then $\Delta=5$ (half zone-width shift, as in Paper~40).
\end{itemize}
Metric: sg4n$_\mathrm{new}$ (under current $\Delta$) and collapse rate per
cell per step in early ($T \leq 2100$) vs.\ late ($T \geq 4200$) windows.

\paragraph{Exp C: Phase aliasing (Anomaly~3).}
$T_\mathrm{switch}\in\{50,100\}$, 10 seeds each, $T=6000$.
Two measurement modes recorded simultaneously:
\begin{itemize}
  \item \textbf{Biased CPS}: sg4n$_A$ and sg4n$_B$ at every 300-step
        checkpoint (reproducing Paper~41's protocol).
  \item \textbf{Epoch-end}: sg4n$_A$ measured at the final step of each
        task~A epoch; sg4n$_B$ at the final step of each task~B epoch
        (post-transient, $t \geq 1500$).
\end{itemize}

\section{Results}

\subsection{Anomaly 1: Copy-Forward Without Waves}

\begin{table}[h]
  \centering
  \caption{Exp~A: sg4n at $T=2000$ and $T=4000$ for each wave condition.
           Ratio $= $ sg4n$(T=4000)/$sg4n$(T=2000)$.
           Field-decay-only prediction: $0.9997^{2000}=0.549$.}
  \label{tab:expa}
  \small
  \begin{tabular}{lcccr}
    \toprule
    Condition & sg4n$(T=2000)$ & sg4n$(T=4000)$ & Ratio & CC/step \\
    \midrule
    C$_\mathrm{ref}$      & 0.397 & 0.077 & 0.195 & 5.71 \\
    C$_\mathrm{wavestop}$ & 0.397 & 0.263 & \textbf{0.664} & 2.84 \\
    C$_\mathrm{nowaves}$  & 0.183 & 0.182 & 0.992 & 0.13 \\
    \midrule
    Field-decay only & --- & --- & 0.549 & --- \\
    \bottomrule
  \end{tabular}
\end{table}

Table~\ref{tab:expa} contains the core result for Anomaly~1.

\paragraph{C$_\mathrm{wavestop}$ decays slower than field-decay-only.}
When waves stop at the commitment peak ($T=2000$, sg4n$=0.397$),
the ratio after 2000 further steps is $0.664$ --- \emph{above} the
field-decay-only prediction of $0.549$.
The 40\% excess ($0.664/0.549 = 1.21$) comes from spontaneous copy-forward
maintenance: cells still collapse occasionally (CC/step$=2.84$, about half
the rate of C$_\mathrm{ref}$) from ageing and instability-noise alone, even
without waves driving viability violations.
Each collapse triggers a birth that propagates fieldM from neighbours,
partially counteracting $\mathrm{FIELD\_DECAY}$.
This explains C$_\mathrm{loc}$ ratio$>1$ from V89: in the location-crystal
condition, collapses at fixed positions seed from neighbourhood fieldM,
maintaining the spatial correlation pattern that sg4n measures.

\paragraph{C$_\mathrm{ref}$ decays faster than field-decay-only.}
Continuous wave input actively overwrites fieldM --- new zone-specific signals
compete with existing structure --- causing a faster decline than pure decay.
After commitment, wave input is a double-edged signal: it refreshes some
zones while introducing noise in others.

\paragraph{C$_\mathrm{nowaves}$ sg4n is stable (scale invariance).}
With WR$=0$, collapse rate is $\approx 0.13$/step (intrinsic noise only).
fieldM decays by $\mathrm{FIELD\_DECAY}$ per step, but sg4n is
\emph{scale-invariant}: both zone-mean distances and within-zone standard
deviations scale proportionally with the fieldM magnitude, leaving their
ratio (sg4n) approximately constant.
The residual sg4n$\approx 0.18$ reflects random spatial correlations
arising from the initial random state; it neither grows nor decays on the
metastability timescale because there is no zone-specific input to drive
differentiation or wash it out.

\subsection{Anomaly 2: C$_\mathrm{half}$ Surge Persists to T=6000}

\begin{table}[h]
  \centering
  \caption{Exp~B: sg4n$_\mathrm{new}$ at key time points.  C$_\mathrm{half}$
           surge peaks between $T=4000$ and $T=6000$, then slowly declines
           but remains $3.1\times$ C$_\mathrm{ref}$ at $T=6000$.}
  \label{tab:expb}
  \small
  \begin{tabular}{lcccc}
    \toprule
    Condition & $T=1500$ (pre-shift) & $T\approx3900$ (peak) & $T=6000$ & $\times$C$_\mathrm{ref}$ \\
    \midrule
    C$_\mathrm{ref}$  & 0.378 & 0.207 & 0.068 & $1.0\times$ \\
    C$_\mathrm{half}$ & 0.378 & 0.558 & \textbf{0.212} & $3.1\times$ \\
    \bottomrule
  \end{tabular}
\end{table}

At $T=6000$, C$_\mathrm{half}$ maintains sg4n$=0.212$ versus C$_\mathrm{ref}=0.068$
(Table~\ref{tab:expb}).
The surge peaks between $T=4000$ and $T=6000$, then gradually declines, but
never approaches the single-task long-run baseline.

\paragraph{Mechanism: sustained boundary-zone collapses.}
After the $\Delta=5$ shift, cells near old zone boundaries occupy mixed
positions under the new labelling (new zone~0 is composed of cells from old
zone~0 and old zone~3).
These cells receive conflicting wave signals --- suppression from one class,
excitation from another --- driving continuous viability violations and
collapses.
Each collapse propagates fieldM via the copy-forward loop, replenishing
the zone structure under the \emph{new} labelling.
This is precisely the mechanism identified in Paper~41 for task switching:
zone-boundary perturbation provides sustained copy-forward refresh.
The difference from Paper~41 is that here the perturbation is spatial
(boundary mismatch rather than task alternation), not temporal.
Both produce the same result: sg4n maintained far above the single-task
long-run baseline because structured collapses continuously refresh fieldM.

\paragraph{Why does the surge exceed the C$_\mathrm{ref}$ peak?}
In C$_\mathrm{ref}$, zone structure peaks at $t^*\approx 2000$ and then
decays as the copy-forward loop slows (fewer viability violations once the
system is committed).
In C$_\mathrm{half}$, the boundary mismatch provides a \emph{continuous}
perturbation that re-enters the system every wave cycle: boundary-zone cells
can never fully commit because the input is internally inconsistent.
This converts what would be a declining metastability trajectory into a
sustained high-structure plateau, at the cost of perfect commitment to
either the old or new zone map.

\subsection{Anomaly 3: Checkpoint Phase Aliasing}

\paragraph{The aliasing mechanism.}
The 300-step checkpoint interval is exactly divisible by both $T_\mathrm{switch}$
values: $300/50=6$ (even multiplier) and $300/100=3$ (odd multiplier).
This has a decisive effect on which task epoch is sampled at each checkpoint.

At $T_\mathrm{switch}=50$: tasks alternate every 50 steps starting with A
(epochs 0--49: A, 50--99: B, \ldots).
The checkpoint at $t=299$ falls in epoch~5 ($=299\div50$, $5\mod 2=1$, B).
The checkpoint at $t=599$ falls in epoch~11 (B), and so on.
Because $300/50=6$ is even, every 300-step increment stays in the same
task-parity: \emph{all 20 checkpoints fall in task~B epochs}.

At $T_\mathrm{switch}=100$: the checkpoint at $t=299$ falls in epoch~2 (A),
at $t=599$ in epoch~5 (B), alternating.
The \emph{last three} checkpoints (at $t=5399, 5699, 5999$) fall in
epochs~53 (B), 56 (A), 59 (B) --- a 2:1 B-weighted average.

\begin{table}[h]
  \centering
  \caption{Exp~C: Biased CPS measurement (reproducing Paper~41 protocol)
           vs.\ epoch-end measurement (unbiased).
           Bold value at $T_\mathrm{switch}=100$ is the phase-aliasing artifact (B$>$A).
           Epoch-end restores sg4n$_A \geq $ sg4n$_B$ for both conditions.}
  \label{tab:expc}
  \small
  \begin{tabular}{rcccc}
    \toprule
    $T_\mathrm{switch}$ & \multicolumn{2}{c}{Biased CPS (last 3)} & \multicolumn{2}{c}{Epoch-end (corrected)} \\
    \cmidrule{2-3} \cmidrule{4-5}
    & sg4n$_A$ & sg4n$_B$ & sg4n$_A$ & sg4n$_B$ \\
    \midrule
    50  & 0.537 & 0.378 & 0.503 & 0.361 \\
    100 & 0.414 & \textbf{0.566} & 0.370 & 0.363 \\
    \bottomrule
  \end{tabular}
\end{table}

\paragraph{Interpretation.}
The biased CPS measurement at $T_\mathrm{switch}=100$ shows sg4n$_B=0.566
>$ sg4n$_A=0.414$ (Table~\ref{tab:expc}).
This is an artifact: the 2:1 B-weighted last-three-checkpoint average
over-samples sg4n$_B$ at its local peak (end of B epochs, when task~B
structure is freshly active).
The epoch-end measurement averages sg4n$_A$ over all task~A epoch ends
and sg4n$_B$ over all task~B epoch ends --- symmetric sampling.
It shows sg4n$_A=0.370 \approx $ sg4n$_B=0.363$: both tasks settle to
essentially equal steady-state sg4n, with a small first-mover advantage
for task~A ($\Delta \approx 0.007$, one standard error).

At $T_\mathrm{switch}=50$, the phase aliasing (all CPS in task~B) does not
cause a reversal because the sawtooth amplitude is small: with only 50-step
epochs, the per-epoch rise and fall in sg4n is minimal, so the epoch-phase
has little effect on the measured value.
The first-mover advantage for task~A is robust regardless of measurement method.

\paragraph{Methodological lesson.}
Checkpoint-based metrics introduce phase bias whenever $T_\mathrm{cp}$ is a
rational multiple of $T_\mathrm{switch}$.
The bias direction and magnitude depend on the parity of the multiplier and
the number of final checkpoints used for averaging.
Epoch-end measurement is the correct protocol for any switching experiment.

\section{Discussion}

\subsection{Unified Picture: The Copy-Forward Loop Is the Maintenance Mechanism}

Anomalies~1 and~2 are both instances of the same principle established in
Paper~41: the copy-forward loop --- collapse~$\rightarrow$ birth~$\rightarrow$
fieldM propagation --- maintains zone structure as long as collapses are
spatially correlated with the existing structure.

In Anomaly~1 (wave ablation), spontaneous collapses from ageing and instability
noise continue after waves stop, providing a low-rate but non-zero copy-forward
signal.
The rate is sufficient to keep the maintenance ratio ($0.664$) above pure
field-decay ($0.549$).
In Anomaly~2 (C$_\mathrm{half}$ surge), the boundary mismatch provides a
\emph{structured} collapse signal at a higher rate, sustaining zone structure
$3.1\times$ above the single-task long-run baseline for at least 4500 steps
beyond the shift.

The unifying principle: \emph{the copy-forward loop requires a perturbation
source, but not necessarily the same perturbation source as the original
learning signal.}
Wave input, task switching (Paper~41), zone-boundary mismatch (Paper~40), and
even ageing noise can all drive the copy-forward loop, with effects proportional
to the structured-collapse rate they generate.

\subsection{The Metastability Half-Life Is Not a Universal Timescale}

Paper~38 measured a single-task half-life of $\approx 800$ steps.
Papers~40--42 collectively show this is a lower bound applicable only under
continuous single-task input with no structured perturbation.
Any persistent source of structured collapses extends this timescale
indefinitely, with the amplification factor set by the ratio of
copy-forward refresh rate to $\mathrm{FIELD\_DECAY}$ rate.
The $0.664$ vs $0.549$ result in Exp~A quantifies the minimum amplification
from intrinsic turnover alone (cc/step $\approx 2.84$ from noise).

\subsection{Completeness of Single-Region Theory}

With these three anomalies resolved, the single-region VCML theory is
internally consistent.
The full picture:
\begin{enumerate}
  \item \textbf{Zone structure formation} (Papers~21--27):
        commitment epoch $t^* \sim 1/(P_c \cdot \mathrm{FA} \cdot \mathrm{WR})$,
        pitchfork mechanism.
  \item \textbf{Metastability} (Papers~28--38):
        half-life $\approx 800$ steps under single-task input;
        crystal conditions modulate the rate.
  \item \textbf{Continual learning} (Papers~39--41):
        zone boundary shifts cause genuine disruption; copy-forward refreshing
        converts task switching from a threat to a maintenance mechanism.
  \item \textbf{Anomaly resolution} (Paper~42):
        spontaneous copy-forward operates below the wave-driven rate;
        boundary mismatch provides sustained structured collapses;
        checkpoint phase aliasing creates measurement artifacts in switching
        experiments.
\end{enumerate}

\section{Conclusion}

Three loose ends from Papers~38--41 are closed:
\begin{enumerate}
  \item \textbf{Spontaneous copy-forward (Anomaly~1):}
        Waves are not necessary for copy-forward maintenance.
        Ageing and instability noise generate $\approx 2.84$ collapses/step
        (half the wave-driven rate), keeping decay ratio at $0.664$ vs
        field-decay-alone prediction of $0.549$.
        C$_\mathrm{loc}$ ratio$>1$ is explained: at fixed positions with sparse
        collapses, copy-forward can slightly amplify zone structure before decay
        dominates.
  \item \textbf{C$_\mathrm{half}$ surge mechanism (Anomaly~2):}
        The half-width boundary shift creates permanent boundary-zone
        ambiguity, driving sustained structured collapses that continuously
        refresh fieldM.
        sg4n$=0.212$ at $T=6000$ (vs.\ C$_\mathrm{ref}=0.068$, $3.1\times$
        amplification) is the long-term steady state of this mechanism,
        identical in kind to Paper~41's task-switching amplification.
  \item \textbf{Phase aliasing artifact (Anomaly~3):}
        $300/T_\mathrm{switch}=6$ (even) causes all checkpoints to fall in
        task~B epochs at $T_\mathrm{switch}=50$.
        $300/T_\mathrm{switch}=3$ (odd) causes 2:1 B-weighting at
        $T_\mathrm{switch}=100$, inflating sg4n$_B$.
        Epoch-end measurement corrects both: sg4n$_A \geq $ sg4n$_B$ in all
        conditions, consistent with task~A's first-mover advantage.
        The apparent reversal in Paper~41 Table~1 was a measurement artifact.
\end{enumerate}

\noindent The single-region VCML theory is now closed.
Future work extends to multi-region substrates
(graph topology, hierarchical relay, sequence learning).

\end{document}
